\section{Preliminaries}
We quickly overview the original integral representation and the ridgelet transform, a mathematical model of depth-2 fully-connected network and its right inverse. Then, we list a few facts in the group representation theory.
In particular, \emph{Schur's lemma} play key roles in the proof of the main results.

\paragraph{Notation.}
For any topological space $X$, $C_c(X)$ denotes the Banach space of all compactly supported continuous functions on $X$.
For any measure space $X$, $L^p(X)$ denotes the Banach space of all $p$-integrable functions on $X$.
$\calS(\RR^d)$ and $\calS'(\RR^d)$ denote the classes of rapidly decreasing functions (or Schwartz test functions) and tempered distributions on $\RR^d$, respectively. 

\subsection{Integral Representation and Ridgelet Transform for Depth-2 Fully-Connected Network}

\begin{definition} For any measurable functions $\sigma:\RR\to\CC$ and $\gamma:\RR^m\times\RR\to\CC$, put
	\begin{align*}
		\lm_\sigma[\gamma](\xx) := \int_{\RR^m\times\RR} \gamma(\aa,b)\sigma(\aa\cdot\xx-b)\dd\aa\dd b, \ \xx \in \RR^m.
	\end{align*}
	We call
	$\lm_\sigma[\gamma]$ an (integral representation of) neural network, and $\gamma$ a parameter distribution.
\end{definition}
The integration over all the hidden parameters $(\aa,b) \in \RR^m\times\RR$ means all the neurons $\{ \xx\mapsto\sigma(\aa\cdot\xx-b) \mid (\aa,b) \in \RR^m\times\RR \}$ are summed (or integrated, to be precise) with weight $\gamma$, hence formally $\lm_\sigma[\gamma]$ is understood as a continuous neural network with a single hidden layer.
We note, however, when $\gamma$ is a finite sum of point measures such as $\gamma_p = \sum_{i=1}^p c_i \delta_{(\aa_i,b_i)}$ (by appropriately extending the class of $\gamma$ to Borel measures), then it can also reproduce a finite width network
\begin{align*}
	\lm_\sigma[\gamma_p](\xx) = \sum_{i=1}^p c_i \sigma(\aa_i\cdot\xx-b_i).
\end{align*}
In other words, the integral representation is a mathematical model of depth-2 network with \emph{any} width (ranging from finite to continuous).

Next, we introduce the ridgelet transform, which is known to be a right-inverse operator to $\lm_\sigma$.
\begin{definition}
	For any measurable functions $\rho:\RR\to\CC$ and $f:\RR^m\to\CC$, put
	\begin{align*}
		\ridge_\rho[f](\aa,b) := \int_{\RR^m} f(\xx) \overline{\rho(\aa\cdot\xx-b)} \dd \xx, \ (\aa,b) \in \RR^m \times \RR. 
	\end{align*}    
	We call $\ridge_\rho$ a ridgelet transform.
\end{definition}
To be precise, it satisfies the following reconstruction formula.
\begin{theorem}[Reconstruction Formula] \label{thm:reconst.ridgelet}
	Suppose $\sigma$ and $\rho$ are a tempered distribution ($\calS'$) and a rapid decreasing function ($\calS$) respectively.
	There exists a bilinear form $\iiprod{\sigma,\rho}$ such that 
	\begin{align*}
		\lm_\sigma \circ \ridge_\rho [f] = \iiprod{\sigma,\rho} f,
	\end{align*}
	for any square integrable function $f \in L^2(\RR^m)$. Further, the bilinear form is given by
		$\iiprod{\sigma,\rho} = \int_{\RR} \sigma^\sharp(\omega) \overline{\rho^\sharp(\omega)} |\omega|^{-m}\dd\omega,$
	where $\sharp$ denotes the 1-dimensional Fourier transform.
\end{theorem}
See \citet[Theorem~6]{Sonoda2021ghost} for the proof.
In particular, according to \citet[Lemma~9]{Sonoda2021ghost}, for any activation function $\sigma$, there always exists $\rho$ satisfying $\iiprod{\sigma,\rho}=1$. 
Here, $\sigma$ being a tempered distribution means that typical activation functions are covered such as ReLU, step function, $\tanh$, gaussian, etc...
We can interpret the reconstruction formula as a universality theorem of continuous neural networks, since for any given data generating function $f$, a network with output weight $\gamma_f = \ridge_\rho[f]$ reproduces $f$ (up to factor $\iiprod{\sigma,\rho}$), i.e. $S[\gamma_f] = f$. In other words, the ridgelet transform indicates how the network parameters should be organized so that the network represents an individual function $f$. 

The original ridgelet transform was discovered by \citet{Murata1996} and \citet{Candes.PhD}. It is recently extended to a few modern networks by the Fourier slice method \citep[see e.g.][]{Sonoda2024fourier}.
In this study, we present a systematic scheme to find the ridgelet transform for a variety of given network architecture based on the group theoretic arguments.


\subsection{Irreducible Representation and Schur's Lemma}
In the main theorem, we use \emph{Schur's lemma}, a fundamental theorem from group representation theory. 
We refer to \citet{Folland2015} for more details on group representation and harmonic analysis on groups.

In this study, we assume group $G$ to be \emph{locally compact}. This is a sufficient condition for having invariant measures. It is not a strong assumption. For example, any finite group, discrete group, compact group, and finite-dimensional Lie group are locally compact, while an \emph{infinite}-dimensional Lie group is \emph{not} locally compact.

Let $\calH$ be a nonzero Hilbert space, and
$\calU(\calH)$ be the group of unitary operators on $\calH$. A \emph{unitary representation} $\pi$ of $G$ on $\calH$ is a group homomorphism that is continuous with respect to the strong operator topology---that is, 
a map $\pi : G \to \calU(\calH)$ satisfying $\pi_{gh} = \pi_g\pi_h$ and $\pi_{g^{-1}}=\pi_g^{-1}$, and for any $\psi \in \calH$, the map $G \ni g \mapsto \pi_g[\psi] \in \calH$ is continuous.

Suppose $\calM$ is a closed subspace of $\calH$. $\calM$ is called an \emph{invariant} subspace when $\pi_g[\calM] \subset \calM$ for all $g \in G$. Particularly, $\pi$ is called \emph{irreducible} when it has only trivial invariant subspaces, namely $\{0\}$ and the whole space $\calH$.
The following theorem is a basic and useful characterization of irreducible representations.
\begin{theorem}[Schur's lemma] \label{thm:schur}
A unitary representation $\pi : G \to \calU(\calH)$ is irreducible iff 
any bounded operator $T$ on $\calH$ that commutes with $\pi$ is always a constant multiple of the identity. In other words, if $\pi_g \circ T=T \circ \pi_g$ for all $g \in G$, then $T=c \id_{\calH}$ for some $c \in \CC$.
\end{theorem}
See %
\citet[Theorem~3.5(a)]{Folland2015} for the proof. %
As we will see in the proof of the main theorem, an irreducible representation (or more generally, a \emph{simple object}) is a standard unit for expressive power. Namely, suppose $X$ is a simple object (such as a simple abelian group, and an irreducible representation), 
and $N$ is a non-trivial sub-object (such as a normal group, and a sub-representation), then we can conclude $N=X$, which means the universality of $N$ in $X$. Schur's lemma restates this in terms of morphism. That is, the commutative property $\pi \circ T = T \circ \pi$ implies $T$ is a homomorphism, and thus it has to be either zero or identity.

As a concrete example of an irreducible representation, we use the following regular representation of the affine group $\aff(m)$ on $L^2(\RR^m)$.
\begin{theorem} \label{thm:aff.irrep}
Let $G := \aff(m) := GL(m) \ltimes \RR^m$ be the affine group acting on $X=\RR^m$ by $(L,\tt) \cdot \xx = L\xx+\tt$, and let $\calH := L^2(\RR^m)$ be the Hilbert space of square-integrable functions.
Let $\pi : \aff(m) \to \calU(L^2(\RR^m))$ be the regular representation of the affine group $\aff(m)$ on $L^2(\RR^m)$, namely $\pi_g[f](\xx) := |\det L|^{-1/2} f( L^{-1}(\xx-\tt) )$ for any $g=(L,\tt) \in G$. Then $\pi$ is irreducible.
\end{theorem}
See \cref{sec:irrep.affine} for the proof.


